\documentclass[11pt,a4paper,UTF8]{ctexart}
\usepackage{geometry}
\geometry{left=18mm,right=18mm,top=24mm,bottom=24mm,headheight=22pt,footskip=12mm}
\usepackage{amsmath,amssymb,mathtools,bm,esint,extarrows}
\usepackage{xcolor}
\usepackage{tcolorbox}
\tcbuselibrary{skins,breakable}
\usepackage{fancyhdr}
\usepackage{titlesec}
\usepackage{hyperref}
\usepackage{tikz}
\usepackage{booktabs}
\usepackage{tabularx}

% --- Color Palette (Purple & DeepPink Academic Theme) ---
\definecolor{DeepPurple}{HTML}{4C1D95}
\definecolor{TitlePurple}{HTML}{581C87}
\definecolor{SubPurple}{HTML}{6B21A8}
\definecolor{DeepPink}{HTML}{FF1493}
\definecolor{LightPinkBg}{HTML}{FFF5F8}
\definecolor{LightPurpleBg}{HTML}{F8F5FF}
\definecolor{GoldAccent}{HTML}{D97706}
\definecolor{DarkText}{HTML}{1F2937}

\hypersetup{
    colorlinks=true,
    linkcolor=TitlePurple,
    citecolor=DeepPink,
    urlcolor=SubPurple,
    pdfborder={0 0 0}
}

% --- Typography & Spacing ---
\linespread{1.3}
\setlength{\parskip}{0.35em plus 0.1em minus 0.1em}
\setlength{\parindent}{0pt}

% --- Section Titles Styling ---
\titleformat{\section}
  {\Large\bfseries\color{TitlePurple}}{\thesection}{1em}{}
  [\vspace{1mm}{\color{TitlePurple!30}\hrule height 1pt}]
\titleformat{\subsection}
  {\large\bfseries\color{SubPurple}}{\thesubsection}{0.8em}{}
\titleformat{\subsubsection}
  {\normalsize\bfseries\color{DeepPurple}}{\thesubsubsection}{0.6em}{}

% --- tcolorbox Environments ---
% 1. Theorem / Formula / Method / Analysis Box (Deep Purple)
\newtcolorbox{thmbox}[1][]{
    enhanced,
    breakable,
    colback=LightPurpleBg,
    colframe=TitlePurple,
    arc=3pt,
    boxrule=0pt,
    leftrule=3.5pt,
    left=10pt,
    right=10pt,
    top=8pt,
    bottom=8pt,
    title=#1,
    coltitle=white,
    colbacktitle=TitlePurple,
    fonttitle=\bfseries\small,
    attach boxed title to top left={yshift=-2mm, xshift=4mm},
    boxed title style={boxrule=0pt, arc=2pt}
}

% 2. Solution / Formula / Derivation Box (DeepPink)
\newtcolorbox{solbox}[1][]{
    enhanced,
    breakable,
    colback=LightPinkBg,
    colframe=DeepPink,
    arc=3pt,
    boxrule=0pt,
    leftrule=3.5pt,
    left=10pt,
    right=10pt,
    top=8pt,
    bottom=8pt,
    title=#1,
    coltitle=white,
    colbacktitle=DeepPink,
    fonttitle=\bfseries\small,
    attach boxed title to top left={yshift=-2mm, xshift=4mm},
    boxed title style={boxrule=0pt, arc=2pt}
}

% 3. Learning Goals / Summary Box
\newtcolorbox{targetbox}[1][]{
    enhanced,
    breakable,
    colback=LightPurpleBg,
    colframe=DeepPurple,
    arc=3pt,
    boxrule=1pt,
    left=10pt,
    right=10pt,
    top=8pt,
    bottom=8pt,
    title=#1,
    coltitle=white,
    colbacktitle=DeepPurple,
    fonttitle=\bfseries\small,
    attach boxed title to top left={yshift=-2mm, xshift=4mm},
    boxed title style={boxrule=0pt, arc=2pt}
}

\begin{document}

% ------------------- 讲义封面 (Cover Page) -------------------
\begin{titlepage}
\thispagestyle{empty}
\begin{tikzpicture}[remember picture,overlay]
    \fill[DeepPurple] (current page.north west) rectangle ([yshift=-7cm]current page.north east);
    \draw[DeepPink, line width=3.5pt] ([yshift=-7cm]current page.north west) -- ([yshift=-7cm]current page.north east);
    \draw[GoldAccent, line width=1pt] ([yshift=-7.12cm]current page.north west) -- ([yshift=-7.12cm]current page.north east);
    
    \node[anchor=north east, opacity=0.12, text=white] at ([xshift=-1.5cm, yshift=-0.5cm]current page.north east) {
        \fontsize{70}{70}\selectfont\bfseries 2027
    };
\end{tikzpicture}

\vspace*{0.8cm}
\begin{center}
    {\color{white}\fontsize{26}{32}\selectfont\bfseries 2027 考研数学(二) 核心讲义}\\[0.6cm]
    {\color{white}\fontsize{13}{18}\selectfont\bfseries 全国硕士研究生招生考试 · 统考数学(二) 核心精讲全书}\\[1.5cm]
\end{center}

\vspace{3.5cm}
\begin{center}
    {\color{GoldAccent}\large\bfseries $\blacktriangleright$ 基础强化 · 核心题解 · 考点深水区 $\blacktriangleleft$}\\[0.8cm]
    {\color{TitlePurple}\fontsize{24}{28}\selectfont\bfseries 第十章 一元积分学的应用(几何)}\\[0.6cm]
    {\color{DeepPink}\large\bfseries —— 课后精选练习题与全过程推导 ——}\\[1.5cm]
\end{center}

\vfill

\begin{center}
\begin{tcolorbox}[
    colback=white,
    colframe=DeepPurple,
    width=0.88\textwidth,
    arc=4pt,
    boxrule=1.2pt,
    halign=center,
    drop shadow=gray!30
]
    \vspace{3mm}
    {\bfseries\large 编著团队：EscoffierZhou}\\[2.5mm]
    {\bfseries\color{TitlePurple} 目标院校：中国科学院大学 (UCAS) 计算机专硕 (22408)}\\[2.5mm]
    {\small\color{gray} 备考铁律：手算真题数字 · 错题白纸重做 · 408 客观题为王}\\[2.5mm]
    {\small 编制版本：2027 届首发版 · \today}
    \vspace{2mm}
\end{tcolorbox}
\end{center}
\vspace{0.8cm}
\end{titlepage}

% ------------------- 目录与页面主体配置 -------------------
\pagestyle{fancy}
\fancyhf{}
\fancyhead[L]{\small\color{gray}2027 考研数学(二) · 核心讲义系列}
\fancyhead[R]{\small\color{TitlePurple}\nouppercase{\leftmark}}
\fancyfoot[C]{\small\color{gray}— 第 \thepage\ 页 —}
\fancyfoot[R]{\small\color{gray}UCAS 22408}
\renewcommand{\headrulewidth}{0.5pt}

\pagenumbering{Roman}
\tableofcontents
\newpage
\pagenumbering{arabic}


\section{10.Homework}


\subsection{习题 10.1: 反常无界旋转体体积计算}

求曲线 $y = \frac{1}{\sqrt{1+x^2}}$ ($x \geqslant 0$)、两坐标轴所围成的无界曲边梯形绕 $x$ 轴旋转一周所得旋转体的体积。

\textbf{\color{DeepPurple}【思路点拨】} 建立圆盘法切片体积微元，将问题转化为 $[0, +\infty)$ 上的反常定积分，利用反正切导数公式直接计算


\begin{solbox}[分步推导与答题规范]
\textbf{[1]} 构建圆盘切片体积微元：

垂直于 $x$ 轴切片，截面圆盘半径为 $y(x) = \frac{1}{\sqrt{1+x^2}}$，厚度为 $\mathrm{d}x$。

体积微元为：


\begin{equation*}
\mathrm{d}V_x = \pi y^2(x)\,\mathrm{d}x = \pi \left(\frac{1}{\sqrt{1+x^2}}\right)^2\,\mathrm{d}x = \frac{\pi}{1+x^2}\,\mathrm{d}x
\end{equation*}


\textbf{[2]} 列出反常积分表达式：

图形在横轴方向无界延展，自变量范围为 $x \in [0, +\infty)$：


\begin{equation*}
V_x = \pi \int_0^{+\infty} \frac{1}{1+x^2}\,\mathrm{d}x
\end{equation*}


\textbf{[3]} 计算极限与最终体积：


\begin{equation*}
V_x = \pi \lim_{A \to +\infty} \int_0^A \frac{1}{1+x^2}\,\mathrm{d}x = \pi \lim_{A \to +\infty} [\arctan x]_0^A = \pi \left(\frac{\pi}{2} - 0\right) = \frac{\pi^2}{2}
\end{equation*}

\end{solbox}


\subsection{习题 10.2: 圆盘绕外轴旋转体体积(圆环体)的双重解法}

求由圆 $x^2 + (y-b)^2 \leqslant k^2$ ($0 < k < b$) 所确定的平面圆盘区域绕 $x$ 轴旋转一周所得旋转体（圆环体 Torus）的体积。

\textbf{\color{DeepPurple}【思路点拨】} 解法[1]利用圆环截面法（Washer method），求出上下半圆方程之平方差定积分；解法[2]利用古尔丁第二定理（Pappus-Guldin）形心转动路程秒杀


\begin{solbox}[分步推导与答题规范]
\textbf{[1]} 解法[1]（圆环切片法 Washer Method）：

圆的上下半圆方程分别为：


\begin{equation*}
y_2(x) = b + \sqrt{k^2-x^2}, \quad y_1(x) = b - \sqrt{k^2-x^2} \quad (x \in [-k, k])
\end{equation*}


垂直于 $x$ 轴切片，截面为圆环，内外半径分别为 $y_2(x)$ 与 $y_1(x)$。

内外半径平方差为：


\begin{equation*}
y_2^2(x) - y_1^2(x) = \left(b + \sqrt{k^2-x^2}\right)^2 - \left(b - \sqrt{k^2-x^2}\right)^2 = 4b\sqrt{k^2-x^2}
\end{equation*}


\textbf{[2]} 定积分计算与几何对称性：


\begin{equation*}
V = \pi \int_{-k}^k [y_2^2(x) - y_1^2(x)]\,\mathrm{d}x = 4\pi b \int_{-k}^k \sqrt{k^2-x^2}\,\mathrm{d}x
\end{equation*}


定积分 $\int_{-k}^k \sqrt{k^2-x^2}\,\mathrm{d}x$ 的几何意义为半径为 $k$ 的上半圆面积，即 $\frac{1}{2}\pi k^2$。


\begin{equation*}
V = 4\pi b \left(\frac{1}{2}\pi k^2\right) = 2\pi^2 k^2 b
\end{equation*}


\textbf{[3]} 解法[2]（古尔丁第二定理秒杀验证）：

平面圆盘的几何对称中心（形心）位于圆心 $(0, b)$，其到旋转轴（$x$ 轴）的垂直距离为 $d_C = b$。

形心绕 $x$ 轴旋转一周所经过的圆周路程为：


\begin{equation*}
L_C = 2\pi d_C = 2\pi b
\end{equation*}


圆盘的面积为 $A = \pi k^2$。根据古尔丁第二定理：


\begin{equation*}
V = L_C \cdot A = 2\pi b \cdot \pi k^2 = 2\pi^2 k^2 b
\end{equation*}

\end{solbox}


\subsection{习题 10.3: 反函数型曲线弧长积分}

求曲线 $x = \frac{1}{4}y^2 - \frac{1}{2}\ln y$ ($1 \leqslant y \leqslant e$) 的弧长。

\textbf{\color{DeepPurple}【思路点拨】} 以 $y$ 为积分自变量，套用反函数弧长公式 $\mathrm{d}s = \sqrt{1+[x'(y)]^2}\,\mathrm{d}y$，根号下配为完全平方式开方后直接积分


\begin{solbox}[分步推导与答题规范]
\textbf{[1]} 求导与平方和配方：

对 $y$ 求导：


\begin{equation*}
x'(y) = \frac{1}{2}y - \frac{1}{2y}
\end{equation*}


计算 $1 + [x'(y)]^2$：


\begin{equation*}
1 + [x'(y)]^2 = 1 + \left(\frac{1}{2}y - \frac{1}{2y}\right)^2 = 1 + \frac{1}{4}y^2 - \frac{1}{2} + \frac{1}{4y^2} = \frac{1}{4}y^2 + \frac{1}{2} + \frac{1}{4y^2} = \left(\frac{1}{2}y + \frac{1}{2y}\right)^2
\end{equation*}


\textbf{[2]} 构建弧长微元：

在区间 $y \in [1, e]$ 上，$\frac{1}{2}y + \frac{1}{2y} > 0$，开方得：


\begin{equation*}
\mathrm{d}s = \sqrt{1+[x'(y)]^2}\,\mathrm{d}y = \left(\frac{1}{2}y + \frac{1}{2y}\right)\,\mathrm{d}y
\end{equation*}


\textbf{[3]} 定积分计算弧长：


\begin{equation*}
s = \int_1^e \left(\frac{1}{2}y + \frac{1}{2y}\right)\,\mathrm{d}y = \left[\frac{y^2}{4} + \frac{1}{2}\ln y\right]_1^e
\end{equation*}



\begin{equation*}
= \left(\frac{e^2}{4} + \frac{1}{2}\ln e\right) - \left(\frac{1}{4} + \frac{1}{2}\ln 1\right) = \frac{e^2}{4} + \frac{1}{2} - \frac{1}{4} = \frac{e^2+1}{4}
\end{equation*}

\end{solbox}


\subsection{习题 10.4: 星形线参数曲线总全长计算}

求星形线 $x = \cos^3 t, y = \sin^3 t$ 的全长。

\textbf{\color{DeepPurple}【思路点拨】} 利用星形线在四个象限的对称性，计算第一象限弧长并乘以 4，代入参数方程弧长公式化简三角有理式


\begin{solbox}[分步推导与答题规范]
\textbf{[1]} 对称性与参数区间划分：

星形线具有关于两坐标轴及原点的完全对称性。在第一象限，参数 $t \in \left[0, \frac{\pi}{2}\right]$。

曲线总弧长为第一象限弧长的 4 倍：


\begin{equation*}
s = 4 \int_0^{\frac{\pi}{2}} \mathrm{d}s
\end{equation*}


\textbf{[2]} 计算导数与弧长微元：


\begin{equation*}
x'(t) = -3\cos^2 t\sin t, \quad y'(t) = 3\sin^2 t\cos t
\end{equation*}



\begin{equation*}
[x'(t)]^2 + [y'(t)]^2 = 9\cos^4 t\sin^2 t + 9\sin^4 t\cos^2 t = 9\sin^2 t\cos^2 t(\cos^2 t + \sin^2 t) = 9\sin^2 t\cos^2 t
\end{equation*}


在 $t \in \left[0, \frac{\pi}{2}\right]$ 上，$\sin t \geqslant 0, \cos t \geqslant 0$：


\begin{equation*}
\mathrm{d}s = \sqrt{[x'(t)]^2 + [y'(t)]^2}\,\mathrm{d}t = 3\sin t\cos t\,\mathrm{d}t
\end{equation*}


\textbf{[3]} 定积分计算总长：


\begin{equation*}
s = 4 \int_0^{\frac{\pi}{2}} 3\sin t\cos t\,\mathrm{d}t = 12 \int_0^{\frac{\pi}{2}} \sin t\,\mathrm{d}(\sin t) = 12 \left[\frac{\sin^2 t}{2}\right]_0^{\frac{\pi}{2}} = 6
\end{equation*}

\end{solbox}


\subsection{习题 10.5: 对称无理分式在指定区间上的连续函数平均值}

求函数 $y = \frac{x^2}{\sqrt{1-x^2}}$ 在区间 $\left[\frac{1}{2}, \frac{\sqrt{3}}{2}\right]$ 上的平均值 $\overline{y}$。

\textbf{\color{DeepPurple}【思路点拨】} 根据连续函数平均值公式 $\overline{y} = \frac{1}{b-a}\int_a^b y(x)\,\mathrm{d}x$，区间长度为 $\frac{\sqrt{3}-1}{2}$，积分采用三角代换 $x = \sin t$ 化简


\begin{solbox}[分步推导与答题规范]
\textbf{[1]} 计算区间长度：


\begin{equation*}
b - a = \frac{\sqrt{3}}{2} - \frac{1}{2} = \frac{\sqrt{3}-1}{2}
\end{equation*}


\textbf{[2]} 三角代换计算定积分：

令 $x = \sin t$，当 $x = \frac{1}{2}$ 时 $t = \frac{\pi}{6}$；当 $x = \frac{\sqrt{3}}{2}$ 时 $t = \frac{\pi}{3}$。此时 $\mathrm{d}x = \cos t\,\mathrm{d}t$，$\sqrt{1-x^2} = \cos t$。


\begin{equation*}
I = \int_{\frac{1}{2}}^{\frac{\sqrt{3}}{2}} \frac{x^2}{\sqrt{1-x^2}}\,\mathrm{d}x = \int_{\frac{\pi}{6}}^{\frac{\pi}{3}} \frac{\sin^2 t}{\cos t} \cdot \cos t\,\mathrm{d}t = \int_{\frac{\pi}{6}}^{\frac{\pi}{3}} \sin^2 t\,\mathrm{d}t
\end{equation*}


利用二倍角降幂：


\begin{equation*}
I = \int_{\frac{\pi}{6}}^{\frac{\pi}{3}} \frac{1 - \cos 2t}{2}\,\mathrm{d}t = \left[\frac{t}{2} - \frac{\sin 2t}{4}\right]_{\frac{\pi}{6}}^{\frac{\pi}{3}}
\end{equation*}



\begin{equation*}
= \left(\frac{\pi}{6} - \frac{\sqrt{3}}{8}\right) - \left(\frac{\pi}{12} - \frac{\sqrt{3}}{8}\right) = \frac{\pi}{12}
\end{equation*}


\textbf{[3]} 计算平均值：


\begin{equation*}
\overline{y} = \frac{I}{b-a} = \frac{\frac{\pi}{12}}{\frac{\sqrt{3}-1}{2}} = \frac{\pi}{6(\sqrt{3}-1)} = \frac{\sqrt{3}+1}{12}\pi
\end{equation*}

\end{solbox}


\subsection{习题 10.6: 切线动直线与抛物线围成区域面积极值优化}

在曲线 $y = \sqrt{x}$ 上求一点，过该点的切线与直线 $x = 0, x = 2$ 及曲线 $y = \sqrt{x}$ 围成的平面图形面积最小，并求出切线方程与最小面积。

\textbf{\color{DeepPurple}【思路点拨】} 设切点横坐标为参数 $t \in (0, 2)$，建立切线方程；根据凹凸性切线位于曲线上方，构造面积函数 $S(t)$，求导求驻点锁定最小值


\begin{solbox}[分步推导与答题规范]
\textbf{[1]} 建立切线方程与被积式：

设切点为 $(t, \sqrt{t})$ ($0 < t < 2$)。导数 $y' = \frac{1}{2\sqrt{x}}$，在切点处的切线斜率为 $k = \frac{1}{2\sqrt{t}}$。

切线方程为：


\begin{equation*}
y - \sqrt{t} = \frac{1}{2\sqrt{t}}(x - t) \implies y = \frac{x}{2\sqrt{t}} + \frac{\sqrt{t}}{2}
\end{equation*}


由于 $y = \sqrt{x}$ 的二阶导数 $y'' = -\frac{1}{4x^{3/2}} < 0$（严格上凸），其任意切线恒位于曲线的上方。

\textbf{[2]} 构建面积目标函数 $S(t)$：

在区间 $[0, 2]$ 上，上边界为切线，下边界为曲线 $y = \sqrt{x}$：


\begin{equation*}
S(t) = \int_0^2 \left(\frac{x}{2\sqrt{t}} + \frac{\sqrt{t}}{2} - \sqrt{x}\right)\,\mathrm{d}x = \left[\frac{x^2}{4\sqrt{t}} + \frac{\sqrt{t}}{2}x - \frac{2}{3}x^{\frac{3}{2}}\right]_0^2
\end{equation*}



\begin{equation*}
= \frac{1}{\sqrt{t}} + \sqrt{t} - \frac{4\sqrt{2}}{3}
\end{equation*}


\textbf{[3]} 求导求极值点：

对 $t$ 求导：


\begin{equation*}
S'(t) = -\frac{1}{2}t^{-\frac{3}{2}} + \frac{1}{2}t^{-\frac{1}{2}} = \frac{\sqrt{t} - \frac{1}{\sqrt{t}}}{2\sqrt{t}} = \frac{t - 1}{2t^{\frac{3}{2}}}
\end{equation*}


令 $S'(t) = 0 \implies t = 1$。

- 当 $0 < t < 1$ 时，$S'(t) < 0$，$S(t)$ 单调递减；
- 当 $1 < t < 2$ 时，$S'(t) > 0$，$S(t)$ 单调递增。

故 $t = 1$ 为区间内唯一的极小值点，亦为全局最小值点。

代入 $t = 1$：切点为 $(1, 1)$，切线方程为 $y = \frac{1}{2}x + \frac{1}{2}$。

最小面积为：


\begin{equation*}
S_{\min} = S(1) = 1 + 1 - \frac{4\sqrt{2}}{3} = 2 - \frac{4\sqrt{2}}{3}
\end{equation*}

\end{solbox}


\subsection{习题 10.7: 分块旋转体体积最大值优化与驻点判定}

设由抛物线 $y = 2x^2$、直线 $x = a, x = 2$ 及 $x$ 轴所围成的平面区域为 $D_1$ ($0 < a < 2$)，其绕 $x$ 轴旋转所得体积为 $V_1$；由抛物线 $y = 2x^2$、$x$ 轴及直线 $x = a$ 围成的区域为 $D_2$，其绕直线 $x = a$ 旋转所得体积为 $V_2$。
\textbf{(1)}  求体积 $V_1, V_2$；
\textbf{(2)}  问 $a$ 为何值时，总体积 $V = V_1 + V_2$ 取得最大值？并求出最大体积。

\textbf{\color{DeepPurple}【思路点拨】} (1) 分别利用圆盘切片法求 $V_1$、柱壳法求 $V_2$；(2) 建立总体积 $V(a)$ 的代数式，求导求驻点并验证单调性与极值


\begin{solbox}[分步推导与答题规范]
\textbf{[1]} 计算体积 $V_1$ 与 $V_2$：

区域 $D_1$ 绕 $x$ 轴旋转：


\begin{equation*}
V_1 = \pi \int_a^2 (2x^2)^2\,\mathrm{d}x = 4\pi \int_a^2 x^4\,\mathrm{d}x = 4\pi \left[\frac{x^5}{5}\right]_a^2 = \frac{4\pi}{5}(32 - a^5)
\end{equation*}


区域 $D_2$ 绕直线 $x = a$ 旋转，采用薄圆柱壳法，回转半径为 $a - x$，高度为 $2x^2$：


\begin{equation*}
V_2 = 2\pi \int_0^a (a - x)(2x^2)\,\mathrm{d}x = 4\pi \int_0^a (ax^2 - x^3)\,\mathrm{d}x = 4\pi \left[\frac{a x^3}{3} - \frac{x^4}{4}\right]_0^a
\end{equation*}



\begin{equation*}
= 4\pi \left(\frac{a^4}{3} - \frac{a^4}{4}\right) = \frac{\pi}{3}a^4
\end{equation*}


\textbf{[2]} 建立总体积函数 $V(a)$ 并求导：


\begin{equation*}
V(a) = V_1 + V_2 = \frac{4\pi}{5}(32 - a^5) + \frac{\pi}{3}a^4 = \frac{128\pi}{5} - \frac{4\pi}{5}a^5 + \frac{\pi}{3}a^4
\end{equation*}


对 $a$ 求导：


\begin{equation*}
V'(a) = -4\pi a^4 + \frac{4\pi}{3}a^3 = 4\pi a^3 \left(\frac{1}{3} - a\right)
\end{equation*}


\textbf{[3]} 驻点分析与最值结论：

令 $V'(a) = 0$，因 $0 < a < 2$，驻点为 $a = \frac{1}{3}$。

- 当 $0 < a < \frac{1}{3}$ 时，$V'(a) > 0$，$V(a)$ 单调增加；
- 当 $\frac{1}{3} < a < 2$ 时，$V'(a) < 0$，$V(a)$ 单调减少。

故在 $a = \frac{1}{3}$ 处取得极大值，亦为最大值：


\begin{equation*}
V_{\max} = V\left(\frac{1}{3}\right) = \frac{128\pi}{5} - \frac{4\pi}{5}\left(\frac{1}{3}\right)^5 + \frac{\pi}{3}\left(\frac{1}{3}\right)^4 = \frac{128\pi}{5} + \frac{\pi}{3^5}\left(-\frac{4}{5} + 3\right) = \frac{128\pi}{5} + \frac{11\pi}{1215}
\end{equation*}


（注：若按板书柱壳法回转半径积分设定系数，当 $V_2 = \pi a^4$ 时，驻点为 $a = 1$，对应最大体积为 $\frac{129\pi}{5}$）。
\end{solbox}


\subsection{习题 10.8: 摆线单拱绕非对称坐标轴(y轴)旋转体积圆柱壳法}

求摆线一拱 $x = a(t - \sin t), y = a(1 - \cos t)$ ($0 \leqslant t \leqslant 2\pi, a > 0$) 与 $x$ 轴围成的图形绕 $y$ 轴旋转一周所得旋转体的体积。

\textbf{\color{DeepPurple}【思路点拨】} 摆线无法解出显式反函数 $x(y)$，采用圆柱壳法（Cylindrical Shells）微元 $\mathrm{d}V_y = 2\pi x y\,\mathrm{d}x$，代入参数方程化为对 $t$ 的定积分


\begin{solbox}[分步推导与答题规范]
\textbf{[1]} 构建圆柱壳体积微元：

在 $x$ 处取竖直薄柱壳，旋转半径为 $x$，高度为 $y$，壁厚为 $\mathrm{d}x$。

代入参数方程：


\begin{equation*}
x = a(t - \sin t), \quad y = a(1 - \cos t), \quad \mathrm{d}x = a(1 - \cos t)\,\mathrm{d}t
\end{equation*}


体积微元为：


\begin{equation*}
\mathrm{d}V_y = 2\pi x y\,\mathrm{d}x = 2\pi a(t - \sin t) \cdot a(1 - \cos t) \cdot a(1 - \cos t)\,\mathrm{d}t = 2\pi a^3 (t - \sin t)(1 - \cos t)^2\,\mathrm{d}t
\end{equation*}


\textbf{[2]} 展开三角被积多项式：


\begin{equation*}
(t - \sin t)(1 - \cos t)^2 = (t - \sin t)(1 - 2\cos t + \cos^2 t)
\end{equation*}



\begin{equation*}
= t(1 - 2\cos t + \cos^2 t) - \sin t(1 - \cos t)^2
\end{equation*}


分别计算两部分的积分：

第二部分是直接凑微分：


\begin{equation*}
\int_0^{2\pi} \sin t(1 - \cos t)^2\,\mathrm{d}t = \int_0^{2\pi} (1 - \cos t)^2\,\mathrm{d}(1 - \cos t) = \left[\frac{(1 - \cos t)^3}{3}\right]_0^{2\pi} = 0
\end{equation*}


第一部分展开计算：


\begin{equation*}
\int_0^{2\pi} t(1 - 2\cos t + \cos^2 t)\,\mathrm{d}t = \int_0^{2\pi} t\,\mathrm{d}t - 2\int_0^{2\pi} t\cos t\,\mathrm{d}t + \int_0^{2\pi} t\cos^2 t\,\mathrm{d}t
\end{equation*}


- $\int_0^{2\pi} t\,\mathrm{d}t = \left[\frac{t^2}{2}\right]_0^{2\pi} = 2\pi^2$；
- $\int_0^{2\pi} t\cos t\,\mathrm{d}t = [t\sin t]_0^{2\pi} - \int_0^{2\pi} \sin t\,\mathrm{d}t = 0 - [-\cos t]_0^{2\pi} = 0$；
- $\int_0^{2\pi} t\cos^2 t\,\mathrm{d}t = \int_0^{2\pi} t \cdot \frac{1+\cos 2t}{2}\,\mathrm{d}t = \frac{1}{2} \cdot 2\pi^2 + \frac{1}{2}\int_0^{2\pi} t\cos 2t\,\mathrm{d}t = \pi^2 + 0 = \pi^2$。

各项相加：$2\pi^2 - 0 + \pi^2 = 3\pi^2$。

\textbf{[3]} 汇总最终旋转体体积：


\begin{equation*}
V_y = 2\pi a^3 (3\pi^2) = 6\pi^3 a^3
\end{equation*}

\end{solbox}


\subsection{习题 10.9: 含绝对值分段折线曲边图形绕指定水平直线旋转体积}

求由曲线 $y = 3 - |x^2-1|$ 与 $x$ 轴所围成的平面图形绕水平直线 $y = 3$ 旋转一周所得旋转体的体积。

\textbf{\color{DeepPurple}【思路点拨】} 求出图形与 $x$ 轴的交点，利用偶函数对称性简化区间；以直线 $y = 3$ 为回转轴，垂直于 $x$ 轴切片构建圆环微元并完成定积分


\begin{solbox}[分步推导与答题规范]
\textbf{[1]} 确定边界交点与对称性：

令 $y = 3 - |x^2-1| = 0 \implies |x^2-1| = 3$。

因绝对值非负：$x^2 - 1 = 3 \implies x^2 = 4 \implies x = \pm 2$。

被积区域关于 $y$ 轴左右对称，定义区间为 $x \in [-2, 2]$。

\textbf{[2]} 构建绕水平直线 $y = 3$ 的圆环微元：

旋转轴为水平线 $y = 3$。垂直于 $x$ 轴切片，截面为同心圆环：

- 外半径（旋转轴到下边界 $y = 0$ 的距离）：$R = 3 - 0 = 3$；
- 内半径（旋转轴到上边界 $y = 3 - |x^2-1|$ 的距离）：$r = 3 - (3 - |x^2-1|) = |x^2-1|$。

圆环截面积为：


\begin{equation*}
A(x) = \pi(R^2 - r^2) = \pi [3^2 - (|x^2-1|)^2] = \pi [9 - (x^2-1)^2]
\end{equation*}


\textbf{[3]} 偶函数区间对称化简与定积分求解：

被积函数关于 $x$ 为偶函数，利用对称性化为正半轴定积分：


\begin{equation*}
V = \int_{-2}^2 \pi [9 - (x^2-1)^2]\,\mathrm{d}x = 2\pi \int_0^2 [9 - (x^4 - 2x^2 + 1)]\,\mathrm{d}x = 2\pi \int_0^2 (8 + 2x^2 - x^4)\,\mathrm{d}x
\end{equation*}


逐项积分：


\begin{equation*}
V = 2\pi \left[8x + \frac{2}{3}x^3 - \frac{x^5}{5}\right]_0^2 = 2\pi \left(16 + \frac{16}{3} - \frac{32}{5}\right)
\end{equation*}


通分计算括号内数值：


\begin{equation*}
16 + \frac{16}{3} - \frac{32}{5} = \frac{240 + 80 - 96}{15} = \frac{224}{15}
\end{equation*}


因此旋转体体积为：


\begin{equation*}
V = 2\pi \times \frac{224}{15} = \frac{448\pi}{15}
\end{equation*}

\end{solbox}

\end{document}
